\chapter{Pre Eclampsia}
\index{Pre Eclampsia}

\medskip
\noindent\textit{\textbf{Under review.} This chapter is an AI-assisted educational draft; it is not a substitute for professional medical advice.}

\section*{Definition and Overview}
a pregnancy complication involving new high blood pressure after 20 weeks with proteinuria or other signs of organ or placental dysfunction.

\section*{Clinical Features}
severe headache, visual disturbance, upper-abdominal pain, sudden swelling, breathlessness, reduced fetal movement, or sometimes no symptoms.

\section*{When to Seek Medical Care}
Seek medical review for new, persistent, worsening, or function-limiting symptoms. Contact your local emergency services for severe breathing difficulty, collapse, sudden neurological change, severe chest or abdominal pain, or any rapidly worsening illness.

\section*{Diagnosis and Investigations}
repeated blood pressure, urine protein, blood tests for liver, kidney and platelet function, fetal growth and wellbeing assessment, and specialist review.

\section*{Management}
urgent obstetric management may include monitoring, blood-pressure treatment, magnesium sulfate for seizure prevention when indicated, corticosteroids for fetal maturation, and planned birth.

\section*{Complications}
Complications depend on the underlying cause and may include progression, effects on other organs, treatment adverse effects, or reduced quality of life. Early assessment can reduce preventable harm.

\section*{Prognosis and Self-Care}
Outcome depends on the cause, severity, complications, and response to treatment. Follow the care plan, keep appointments, avoid known triggers, and seek help if symptoms worsen.
\section*{Safety-netting}
Seek urgent help for severe or rapidly worsening symptoms, collapse, breathing difficulty, severe pain, confusion, dehydration, uncontrolled bleeding, or any immediate danger.
